\uuid{Ses5}
\exo7id{3466}
\titre{exo7 3466}
\auteur{quercia}
\organisation{exo7}
\datecreate{2010-03-10}
\isIndication{false}
\isCorrection{true}
\chapitre{Déterminant, système linéaire}
\sousChapitre{Système linéaire, rang}
\cpge{mpsi-mp2i.matrices-applications-lineaires.b, pcsi.matrices-applications-lineaires.b, ptsi.matrices-applications-lineaires.b, mpsi-mp2i.calcul-matriciel-systemes.e, pcsi.determinants.c, ptsi.determinants.a, bcpst1.matrices, bcpst2.applications-lineaires-matrices.b}
\module{Algèbre}
\niveau{L2}
\difficulte{}

\contenu{
\question{
Soit la matrice carrée d'ordre~$n$, $I_p$ ($p\le n$), telle que le $i$-ème
terme diagonal vaut 1 si $i$ est compris entre $p$ et~$n$, tous les autres
coefficients étant nuls.
Quelle sont les conditions sur~$A$ (matrice carrée d'ordre~$n$) pour
qu'il existe~$B$ telle que $AB=I_p$~?
}
\reponse{
Les colonnes de~$A$ engendrent les $n-p$ derniers vecteurs de la
base canonique.
}
}
